%---------- 第二章 相关理论与技术基础 ----------
\section{相关理论与技术基础}

\subsection{卷积神经网络}

\quad 卷积神经网络通过卷积层、池化层与全连接层的堆叠实现对图像特征的逐层提取。
设输入特征图为 $X \in \mathbb{R}^{H \times W \times C}$，卷积核为
$K \in \mathbb{R}^{k \times k \times C}$，则卷积层的输出可表示为
\begin{equation}
  Y_{i,j} = \sigma\left( \sum_{m=1}^{k} \sum_{n=1}^{k} \sum_{c=1}^{C}
    K_{m,n,c} \, X_{i+m-1,\,j+n-1,\,c} + b \right),
  \label{eq-conv}
\end{equation}
其中 $b$ 为偏置项，$\sigma(\cdot)$ 为非线性激活函数。式~\eqref{eq-conv} 表明，
卷积运算通过局部连接与权值共享有效降低了模型参数量。

\subsection{注意力机制}

\quad 注意力机制通过为特征图的不同空间位置分配不同权重，引导网络聚焦于对识别结果
贡献更大的区域。给定查询向量 $\boldsymbol{q}$ 与键值对集合
$\{(\boldsymbol{k}_i, \boldsymbol{v}_i)\}_{i=1}^{N}$，注意力权重定义为
\begin{equation}
  \alpha_i = \frac{\exp\left( \boldsymbol{q}^{\mathsf{T}} \boldsymbol{k}_i
    \middle/ \sqrt{d} \right)}
    {\sum_{j=1}^{N} \exp\left( \boldsymbol{q}^{\mathsf{T}} \boldsymbol{k}_j
    \middle/ \sqrt{d} \right)},
  \qquad
  \mathrm{Attn}(\boldsymbol{q}) = \sum_{i=1}^{N} \alpha_i \boldsymbol{v}_i,
\end{equation}
其中 $d$ 为缩放因子。常用的评价指标包括准确率、召回率与平均精度均值（mAP）。

\subsection{常用数据集}

\quad 本文实验所用公开数据集的基本情况如表~\ref{tab-datasets} 所示。

\begin{table}[htbp]
  \centering
  \caption{常用图像识别数据集概况}
  \label{tab-datasets}
  \begin{tabular}{cccc}
    \toprule
    数据集 & 图像数量 & 类别数 & 图像尺寸 \\
    \midrule
    MNIST     & 70\,000  & 10  & $28 \times 28$ \\
    CIFAR-10  & 60\,000  & 10  & $32 \times 32$ \\
    ImageNet  & $1.4 \times 10^{7}$ & 1\,000 & 不定 \\
    \bottomrule
  \end{tabular}
\end{table}
